We deploy LSE inside a symbolic execution-aware virtual machine (the symbolic VM) that encapsulates the ``entire universe'' of the application, including the framework and even the language interpreter and operating system, enabling integration testing of the entire application stack.
%
The test-mode deployment then requires just the push of a button for the system to execute the layered parameterized tests, generate coverage statistics, and highlight any failing test cases.

This deployment model leverages the properties of PaaS in several ways:
%
(1)~By hiding the testing VMs behind a service interface, the PaaS system can faithfully reproduce the exact environment of production VMs inside the testing VMs without exposing its internals.
%
(2)~The testing task can be transparently scaled across multiple VMs by using parallel symbolic execution~\cite{cloud9}.
%
(3)~Since the application uses standard interfaces for accessing the PaaS components (storage, networking, etc.), the provider is able to substitute production-optimized implementations with testing-optimized stubs that offer a simplified behavior that is better suited to automated program analysis.

From the perspective of the PaaS provider, the test runner service consists of a set of symbolic VMs, operated separately from the production infrastructure.  When an application is deployed in test mode, one of the symbolic VMs is allocated for testing: the application code and tests are copied to the guest, and the LSE algorithm is invoked.

\paragraph{Architecture}

\begin{figure}
  \centering
  \includegraphics[width=3.1in]{paas/figures/symbolic-vm}
  \caption{The PaaS test runner service.}
  \label{fig:fse}
\end{figure}

Figure~\ref{fig:fse} illustrates the architecture of the symbolic VM environment.  Inside the VM~\cI, all application components are symbolically executed in the same low-level representation (e.g., x86 machine code or LLVM~\cite{llvm}).  The components execute inside their own vanilla interpreters~\cII.
%
The test framework~\cIII plays two roles: it implements (1) the APIs for LPTs and onion objects that developers use to write the testing code, and (2) the LSE algorithm that guides the test case generation.

\paragraph{Prototype}

We implemented a prototype of the symbolic VM that tests Python-based Google App Engine PaaS applications and is built on top of the S2E symbolic virtual machine.

In our implementation, the symbolic execution engine and the LSE logic live at different levels in the symbolic VM stack.
%
The symbolic execution engine operates with low-level abstractions such as memory bytes. It resides on the host level, as an S2E plugin that exposes the core symbolic execution primitives to the guest as S2E system calls, e.g., to allow marking memory buffers as symbolic.
%
The LSE algorithm operates on web application state (e.g., by accessing the onion objects), and is implemented in the guest as a native Python extension module.  We implemented LSE on top of WebTest~\cite{py-webtest}, a popular fixture library for testing Python web applications.
%
The resulting system is extensible to other languages with limited engineering effort: since the symbolic execution logic is provided at the host level, only the test framework component needs to be implemented in the cloud app language.

Early experiences with the prototype are encouraging: for the full application stack of a simple cloud app, our prototype generates a test case every few seconds.

%%% Local Variables: 
%%% mode: latex
%%% eval: (visual-line-mode)
%%% fill-column: 1000000
%%% TeX-master: "main"
%%% End:
